\section{Mathematical Induction}
\vspace{-0.5em}

\subsection{Sum of Odd Integers}
\textbf{Prove:}
$1+3+\cdots+(2n-1)=n^2$, $n\geq1$.

\textbf{Base step:} $n=1$: $1=1^2$.

\textbf{Induction Hypothesis:} \\ 
Assume that the statement holds for \\ $n=k, k \in \mathbb{Z}^{+}$. \\
$1+3+\cdots+(2k-1)=k^2$.

\textbf{Inductive step:}
\[
\begin{aligned}
1+3+\cdots+(2k+1)
&=k^2+(2k+1)\\
&=(k+1)^2.
\end{aligned}
\]
so that $P(k+1)$ is true if $P(k)$ is true. \\
Thus, the statement holds for all $n \in \mathbb{Z}^{+}$. $\square$

\subsection{Inequality: $n<2^n$}
\textbf{Prove:} $n<2^n$, $n\geq1$.

\textbf{Base step:} $n=1$: $1<2$. \\
Thus the statement holds for $n=1$

\textbf{Inductive Hypothesis:} \\ 
Assume that the statement holds for \\ 
$n=k, k \in \mathbb{Z}^{+}, k\geq1$ \\
$k<2^k$.

\textbf{Inductive step:} Since $k\geq1$,
\[
k+1\leq2k<2\cdot2^k=2^{k+1}.
\]
so that $P(k+1)$ is true if $P(k)$ is true. \\
Thus, $n<2^{n}$ is true for all $n\in\mathbb{Z}^{+}$. $\square$

\subsection{Divisibility by 3}
\textbf{Prove:} $3\mid(n^3-n)$, $n\geq1$.

\textbf{Base step:} $n=1$: $1^3-1=0$.

\textbf{Inductive Hypothesis:} \\ 
Assume that the statement holds for \\
$n = k,k\in\mathbb{Z}^{+}$
$3\mid(k^3-k)$.

\textbf{Inductive step:}
\[
\begin{aligned}
(k+1)^3-(k+1)
&=k^3+3k^2+2k\\
&=(k^3-k)+3k(k+1).
\end{aligned}
\]
By the inductive step, we conclude that
both terms are divisible by 3.
Thus, $3\mid(n^3-n)$ for all $n\in\mathbb{Z}^{+}$.
$\square$

\subsection{Sum of Squares}
\textbf{Prove:}
\[
\sum_{i=1}^{n}i^2
=\frac{n(n+1)(2n+1)}{6}.
\]

\textbf{Base step:} $n=1$:
$1=\frac{1(2)(3)}{6}$.\\
Thus the statement holds for $n=1$.\\
\textbf{Inductive Hypothesis:} \\ 
Assume the statements holds for \\ $n=k,k\in\mathbb{Z}^{+}$.
\[
\sum_{i=1}^{k}i^2
=\frac{k(k+1)(2k+1)}{6}.
\]

\textbf{Inductive step:}
\[
\begin{aligned}
\sum_{i=1}^{k+1}i^2
&=\frac{k(k+1)(2k+1)}{6}+(k+1)^2\\
&=\frac{(k+1)[k(2k+1)+6(k+1)]}{6}\\
&=\frac{(k+1)(2k^2+7k+6)}{6}\\
&=\frac{(k+1)(k+2)(2k+3)}{6}.
\end{aligned}
\]
This is the formula for $n=k+1$.\\
By mathematical induction, the formula is true for all $n\in\mathbb{Z}^{+}$. $\square$

\subsection{Factorial Inequality}
\textbf{Prove:} $2^n<n!$, $n\geq4$.

\textbf{Base step:} $n=4$: $16<24$.\\
Thus the statement holds for $n=4$
\textbf{Inductive Hypothesis:} 
Assume the statement holds for $n=k,k\in\mathbb{Z}^{+},k\geq4$. \\
$2^k<k!$

\textbf{Inductive step:} Since $k+1>2$,
\[
2^{k+1}=2\cdot2^k
<2k!<(k+1)k!=(k+1)!.
\]
Thus, $P(k+1)$ is true if $P(k)$ is true, $2^n < n!$ for all $n\in\mathbb{Z}^{+}$ where $n\geq4$. $\square$

\subsection{Geometric Sum}
\textbf{Prove:}
\[
1+2+2^2+\cdots+2^n=2^{n+1}-1,
\quad n\geq0.
\]

\textbf{Base step:} $n=0$: $1=2^1-1$.
Thus, the statement holds for $n=0$ \\
\textbf{Inductive Hypothesis:} \\
Assume that the statement holds for \\ $n=k,k\in\mathbb{Z}_{\geq0}$ \\
$\sum_{i=0}^{k}2^i=2^{k+1}-1$.

\textbf{Inductive Step:}
\[
\begin{aligned}
\sum_{i=0}^{k+1}2^i
&=(2^{k+1}-1)+2^{k+1}\\
&=2\cdot2^{k+1}-1\\
&=2^{k+2}-1.
\end{aligned}
\]
By mathematical induction, $P(n)$ is true for all $n\geq0$. $\square$

\subsection{Divisibility by 57}
\textbf{Prove:}
$57\mid(7^{n+2}+8^{2n+1})$, $n\geq0$.

\textbf{Base step:} $n=0$:
\[
7^2+8^1=49+8=57.
\]

\textbf{Inductive Hypothesis:} \\
Assume the statement holds for \\ $n=k,k\in\mathbb{Z},k\geq0$ \\
$57\mid(7^{k+2}+8^{2k+1})$.

\textbf{Inductive step:}
\[
\begin{aligned}
&7^{k+3}+8^{2k+3}\\
&=7\cdot7^{k+2}+64\cdot8^{2k+1}\\
&=7(7^{k+2}+8^{2k+1})
  +57\cdot8^{2k+1}.
\end{aligned}
\]
By the inductive step, we conclude that first term is divisible by $57$,
and the second term is also divisible by $57$.
Thus, $7^{k+2}+8^{2k+1}$ is divisible by $57$ for all $n\in\mathbb{Z}$ where $n\geq0$